\section{Data Processing}

The data-processing layer serves two different modeling burdens. For the evaluation model, the raw supplier indicators must be converted into dimensionless comparison variables so that cost, reliability, fulfillment, and loss behavior can be compared within one score system. For the planning model, the demand sequence, available capacity, and effective arrival quantity must be represented in a form that can be directly inserted into feasibility constraints. Thus, the data section is not a decorative preface; it is the bridge that lets the later evaluation and planning layers share one consistent variable system.

The first step is indicator cleaning and orientation adjustment. Indicators such as reliability or fulfillment rate are benefit-oriented and can enter the score model after normalization, whereas indicators such as unit cost or transport loss are cost-oriented and must be transformed so that a larger normalized value still means better comprehensive performance. This step prevents the later evaluation result from being distorted by inconsistent indicator direction. If the paper uses entropy weighting, AHP, or a hybrid score, the normalization method should be stated here rather than left implicit in the model section.

The second step is candidate-set preparation for planning. After the supplier indicators are normalized, the paper should extract a reduced candidate pool instead of forcing every supplier directly into the decision model. This is where the evaluation route pays off: the planning model can focus on the most credible suppliers, which reduces combinational complexity and also makes the final recommendation easier to interpret. In a complete paper, this step is often supported by a normalized-indicator table, a score-distribution figure, or a retained-candidate summary table.

The third step is demand and feasibility preprocessing. The demand sequence must be aligned with the decision periods, and the effective available quantity should already account for expected loss or other delivery adjustments before the planning constraints are written. If scenario analysis is planned, the disturbed demand, loss, or capacity settings should also be defined at this stage so that the validation section later reads like an extension of the same data logic rather than as an unrelated appendix experiment.

Therefore, the real role of the data section is to stabilize the evidence chain. By the end of this layer, the reader should know which indicators support supplier evaluation, which reduced candidate set enters the planning model, and which scenario parameters are reserved for later robustness checks. Once these pieces are visible, the transition into model establishment becomes natural instead of abrupt.
